\documentclass[10pt,aspectratio=169]{beamer}
% Preámbulo modular para presentaciones Beamer (Modelación Estadística)
% Diseño y paleta institucional del Tecnológico de Monterrey

\documentclass[aspectratio=169,xcolor={dvipsnames,table}]{beamer}

\usepackage[utf8]{inputenc}
\usepackage[spanish,mexico]{babel}
\usepackage{amsmath,amssymb,amsfonts,mathtools}
\usepackage{graphicx}
\usepackage{listings}
\usepackage{booktabs}
\usepackage{tikz}
\usetikzlibrary{matrix,arrows,decorations.pathmorphing,shapes.geometric,positioning}

% Paleta Institucional Tecnológico de Monterrey
\definecolor{TecAzul}{HTML}{0039A6}       % Azul TEC: color primario institucional
\definecolor{TecAzulOscuro}{HTML}{1A2E51} % Azul marino oscuro
\definecolor{TecRojo}{HTML}{EC2661}       % Rojo de acento (paleta miTec)
\definecolor{TecGris}{HTML}{646464}       % Gris neutro para comentarios y subtítulos
\definecolor{TecFondoBloque}{HTML}{F4F6F9}% Fondo suave para bloques

% Configuración de Tema Beamer
\usetheme{Boadilla}
\usecolortheme[named=TecAzulOscuro]{structure}
\setbeamercolor{alerted text}{fg=TecRojo}
\setbeamercolor{palette primary}{bg=TecAzulOscuro,fg=white}
\setbeamercolor{palette secondary}{bg=TecAzul,fg=white}
\setbeamercolor{palette tertiary}{bg=TecAzulOscuro!80!black,fg=white}
\setbeamercolor{title}{fg=TecAzulOscuro,bg=TecFondoBloque}
\setbeamercolor{frametitle}{fg=TecAzulOscuro,bg=TecFondoBloque}

% Bloques personalizados elegantes
\setbeamercolor{block title}{bg=TecAzulOscuro,fg=white}
\setbeamercolor{block body}{bg=TecFondoBloque,fg=black}
\setbeamercolor{block title alerted}{bg=TecRojo,fg=white}
\setbeamercolor{block body alerted}{bg=TecRojo!8,fg=black}
\setbeamercolor{block title example}{bg=TecAzul,fg=white}
\setbeamercolor{block body example}{bg=TecAzul!8,fg=black}

% Redondear bordes y añadir sombra sutil
\setbeamertemplate{blocks}[rounded][shadow=true]
\setbeamertemplate{navigation symbols}{}

% Pie de página limpio con número de diapositiva
\setbeamertemplate{footline}{
  \leavevmode%
  \hbox{%
  \begin{beamercolorbox}[wd=.7\paperwidth,ht=2.25ex,dp=1ex,left]{author in head/foot}%
    \hspace*{2ex}\insertshortauthor~~(\insertshortinstitute) \hfill \insertshorttitle
  \end{beamercolorbox}%
  \begin{beamercolorbox}[wd=.3\paperwidth,ht=2.25ex,dp=1ex,right]{date in head/foot}%
    \insertshortdate{}\hspace*{2em}
    \insertframenumber{} / \inserttotalframenumber\hspace*{2ex}
  \end{beamercolorbox}}%
  \vskip0pt%
}

% Configuración de Listings de Python para visualización pedagógica de código
\lstset{
    language=Python,
    basicstyle=\ttfamily\footnotesize,
    keywordstyle=\color{TecAzulOscuro}\bfseries,
    stringstyle=\color{TecRojo},
    commentstyle=\color{TecGris}\itshape,
    showstringspaces=false,
    breaklines=true,
    frame=lines,
    rulecolor=\color{TecAzulOscuro!30},
    backgroundcolor=\color{TecFondoBloque},
    aboveskip=0.5em,
    belowskip=0.5em,
    literate={á}{{\'a}}1 {é}{{\'e}}1 {í}{{\'i}}1 {ó}{{\'o}}1 {ú}{{\'u}}1
             {Á}{{\'A}}1 {É}{{\'E}}1 {Í}{{\'I}}1 {Ó}{{\'O}}1 {Ú}{{\'U}}1
             {ñ}{{\~n}}1 {Ñ}{{\~N}}1 {¿}{{?`}}1 {¡}{{!`}}1
}

% Metadatos institucionales por defecto
\title[Modelación Estadística]{Modelación Estadística para Ciencia de Datos e Ingeniería}
\author[J. Castillo Colmenares]{Juliho David Castillo Colmenares}
\institute[Tec de Monterrey]{Tecnológico de Monterrey \\ Escuela de Ingeniería y Ciencias}
\date{\today}
% Comandos y macros estadísticas compartidas para las presentaciones Beamer (Metropolis)

% Conjuntos numéricos básicos
\newcommand{\R}{\mathbb{R}}
\newcommand{\N}{\mathbb{N}}
\newcommand{\Q}{\mathbb{Q}}
\newcommand{\Z}{\mathbb{Z}}
\newcommand{\set}[1]{\left\{ #1 \right\}}
\newcommand{\abs}[1]{\left|#1\right|}
\newcommand{\norm}[1]{\left\| #1 \right\|}

% Comandos estadísticos y probabilísticos del curso
\newcommand{\Var}{\operatorname{Var}}
\newcommand{\cov}{\operatorname{Cov}}
\newcommand{\corr}{\rho}
\newcommand{\comb}[2]{\begin{pmatrix} #1 \\ #2 \end{pmatrix}}
\newcommand{\s}{\sigma}
\newcommand{\particion}{\mathfrak{P}}
\newcommand{\card}[1]{\#\left( #1 \right)}
\newcommand{\seq}[3]{\set{#1}_{#2}^{#3}}
\newcommand{\E}{\mathbb{E}}
\newcommand{\Prob}{\mathbb{P}}

% Entornos pedagógicos y cajas en español académico natural
\newenvironment{discusion}{%
  \begin{alertblock}{Pregunta de Discusión}%
}{%
  \end{alertblock}%
}

\newenvironment{reflexion}{%
  \begin{alertblock}{Reflexión para el Aula}%
}{%
  \end{alertblock}%
}

\newenvironment{paradoja}{%
  \begin{alertblock}{Paradoja de Exploración}%
}{%
  \end{alertblock}%
}

\newenvironment{motivacion}{%
  \begin{exampleblock}{Motivación: Ciencia de Datos en la Práctica}%
}{%
  \end{exampleblock}%
}

\newenvironment{retoclase}{%
  \begin{block}{Reto Rápido en Equipo}%
}{%
  \end{block}%
}

\title[Prueba $t$ para una media]{Sección 07.04: Prueba $t$ para una media}\subtitle{Varianza poblacional desconocida}
\author[J. Castillo Colmenares]{Juliho Castillo Colmenares}\institute{Tecnológico de Monterrey}\date{}
\begin{document}
\begin{frame}[plain]\titlepage\end{frame}
\begin{frame}{Hoja de ruta}\small\begin{enumerate}\item De $Z$ a $t$.\item Estadística y distribución.\item Aplicación a una media.\end{enumerate}\end{frame}
\begin{frame}{Fuente y alcance}\small La sección estudia pruebas para una media cuando $\sigma$ es desconocida y debe estimarse mediante $S$.\begin{block}{Pregunta guía}¿Cómo cambia la prueba al estimar la desviación estándar?\end{block}\end{frame}
\begin{frame}{Por qué aparece $t$}\small Sustituir $\sigma$ por $S$ introduce incertidumbre adicional. Por ello el estadístico ya no sigue $N(0,1)$, sino una distribución $t$ de Student.\end{frame}
\begin{frame}{Hipótesis}\small\[H_0:\mu=\mu_0\qquad\text{frente a}\qquad H_a:\mu\neq\mu_0,\ \mu<\mu_0\ \text{o}\ \mu>\mu_0.\]\end{frame}
\begin{frame}{Estadística de prueba}\small\[T=\frac{\bar X-\mu_0}{S/\sqrt n}.\]\begin{block}{Distribución bajo $H_0$} $T\sim t_{n-1}$ cuando la población es normal.\end{block}\end{frame}
\begin{frame}{Grados de libertad}\small La estimación de $S$ consume una cantidad de información; por eso la distribución usa $n-1$ grados de libertad.\end{frame}
\begin{frame}{Regiones de rechazo}\small\begin{itemize}\item Cola derecha: $T>t_{\alpha,n-1}$.\item Cola izquierda: $T<-t_{\alpha,n-1}$.\item Dos colas: $|T|>t_{\alpha/2,n-1}$.\end{itemize}\end{frame}
\begin{frame}{Intervalo asociado}\small El intervalo bilateral correspondiente es\[\bar X\pm t_{\alpha/2,n-1}\frac{S}{\sqrt n}.\]\end{frame}
\begin{frame}{Procedimiento I}\small\begin{enumerate}\item Verificar independencia y normalidad aproximada.\item Formular hipótesis.\item Fijar $\alpha$.\end{enumerate}\end{frame}
\begin{frame}{Procedimiento II}\small\begin{enumerate}\setcounter{enumi}{3}\item Calcular $T$.\item Usar $t_{n-1}$ para valor crítico o valor $p$.\item Redactar la conclusión.\end{enumerate}\end{frame}
\begin{frame}{Supuestos}\small La prueba requiere una muestra aleatoria e independiente. La normalidad es especialmente importante con tamaños pequeños; para tamaños grandes se admite una aproximación por el teorema central del límite.\end{frame}
\begin{frame}{Aplicación: datos}\small Un fabricante afirma una duración media de $500$ horas. Una muestra de $n=16$ baterías tiene $\bar x=485$ y $s=32$.\begin{block}{Alternativa} Se desea saber si la media real es menor.\end{block}\end{frame}
\begin{frame}{Aplicación: hipótesis}\small\[H_0:\mu=500,\qquad H_a:\mu<500.\]El nivel de significación es $\alpha=0.05$ y los grados de libertad son $15$.\end{frame}
\begin{frame}{Aplicación: estadística}\small\[T=\frac{485-500}{32/\sqrt{16}}=-1.875.\]\begin{block}{Lectura} El signo negativo concuerda con una media muestral menor que la referencia.\end{block}\end{frame}
\begin{frame}{Aplicación: decisión}\small Se compara $T$ con el cuantil inferior de $t_{15}$, o se calcula el valor $p$ de cola izquierda. La decisión depende de si el estadístico cae en la región crítica.\end{frame}
\begin{frame}{Aplicación: intervalo}\small El intervalo bilateral utiliza $t_{0.025,15}$ y cuantifica qué medias son compatibles con los datos. Si $500$ queda por encima del extremo superior, la evidencia favorece una media menor.\end{frame}
\begin{frame}{Interpretación}\small Rechazar $H_0$ significa evidencia de una duración media menor a $500$ horas al nivel elegido; no significa que cada batería dure menos.\end{frame}
\begin{frame}{Relación con $Z$}\small Cuando $n$ crece, $t_{n-1}$ se aproxima a $N(0,1)$. En muestras pequeñas, los cuantiles $t$ son más alejados y la prueba es más conservadora.\end{frame}
\begin{frame}{Aplicación de reporte}\small\begin{block}{Plantilla} ``Con $T=\cdots$, $gl=\cdots$ y $p=\cdots$, al nivel $\alpha=\cdots$ [rechazamos/no rechazamos] $H_0$; la evidencia [apoya/no apoya] $\mu<\mu_0$.''\end{block}\end{frame}
\begin{frame}{Síntesis}\small\begin{itemize}\item Estimar $\sigma$ reemplaza $Z$ por $t$.\item La estadística usa $S/\sqrt n$.\item Los grados de libertad son $n-1$.\end{itemize}\end{frame}
\begin{frame}{Conexión con el capítulo}\small La prueba para una media es el modelo base para comparar medias, proporciones y varianzas en las secciones siguientes.\end{frame}
\end{document}
